\scanpage{98}
On sait que cette dernière est équivalente à la cat. ponctuelle. Donc $\mathrm{DR}_{\unsure{\mathrm{sp}}}(X,k)_0$ est équivalente à la cat. ponctuelle.

Si maintenant on travaille \unsure{des|dans} $\mathrm{DR}_{\unsure{\mathrm{spufl}}}$ \deleted{(\unsure{noté} $\mathrm{DR}_{\unsure{\mathrm{spfl}}}$)} i.e. avec des $A^{**}$ \inserted{u.} flasques, ($A^{ij}$ \inserted{u.} flasque si $\unsure{i<j}$) donc $\Phi$ est \inserted{u.} flasque.
\deleted{(\unsure{pour}\ \unsure{flasque})}.

Supposons que ceci implique que les $\Gamma_k^i\Phi$ $(i>0)$ sont uni. flasques, alors $\mathrm{DR}(\Phi,T)$ est aussi flasques, \deleted{et} et même univ. flasque, et on voit comme dessus que la cat. dérivée est équiv. à celle construite avec \deleted{la cat. des} $\mathrm{DR}(\Phi,T)$, avec $\Phi$ univ. flasque et $\Phi/k.T$ plat, i.e. avec les suites exactes
\diagram{\begin{tikzcd}
0\arrow[r]&k\arrow[r]&\Phi\arrow[r]&\Psi\arrow[r]&0
\end{tikzcd}}
avec $\Phi$ univ. flasque et $\Psi$ $k$-plat ; or cette cat. dérivée est encore la catégorie ponctuelle. On trouve donc

\uline{Prop} La cat. $\mathrm{DR}_{\unsure{\mathrm{sp}}}(X,k)_0$ est équivalente à la cat. ponctuelle. Il en est de même de la cat. $\mathrm{DR}_{\unsure{\mathrm{spufl}}}(X,k)_0$, si on suppose que pour tout \inserted{$k$-Module} $\Phi$ plat et univ. flasque, les $\Gamma^i\Phi$ $(i>0)$ sont aussi univ. flasques.

\margin{97}
